\chapter{宋辽夏金元卷读者索引}
\label{app:index}

本索引只链接已经在连续叙事首次呈现的事件锚点；它按行动者、问题与材料入口提供读路，而不把二千二百条账本重印成姓名表。日期以并行政权的原有纪年和账本换算为准。

\section{政权与区域入口}
\subsection{北宋}
\begin{description}
  \item[960年] \eventref{NS-0960-KAIFENG-CAPITAL}{宋廷承用开封为东京和中央行政中心}（宋）。
  \item[960年（宋建隆元年；公元换算据《宋史》本纪纪年）] \eventref{SYD-960-NS-001}{八月戊辰朔，御崇元殿，行入閣儀}（宋）。
  \item[960年（宋建隆元年；公元换算据《宋史》本纪纪年）] \eventref{SYD-960-NS-168}{丙辰，南唐主李景、吳越王錢俶遣使以御服、錦綺、金帛來賀}（宋）。
  \item[961年前后] \eventref{NS-0961-MILITARY-POWER}{太祖以宴饮、任官和调动方式削弱禁军宿将的独立兵权}（宋）。
\end{description}
\subsection{辽／契丹}
\begin{description}
  \item[907年] \eventref{LIAO-0907-EVENT-026}{耶律阿保机取得汗位，开始重组部族权力}（契丹）。
  \item[916年] \eventref{LIAO-0916-EVENT-027}{阿保机称帝并建国号契丹}（契丹）。
  \item[926年] \eventref{LIAO-0926-EVENT-028}{辽军攻灭渤海，设置东丹等安排}（辽与渤海）。
  \item[927年] \eventref{LIAO-0927-EVENT-029}{耶律德光即位为辽太宗}（辽）。
\end{description}
\subsection{西夏／党项}
\begin{description}
  \item[982年] \eventref{XIA-0982-EVENT-049}{李继迁脱离宋朝羁縻控制，在西北边地集结力量}（党项李氏与宋）。
  \item[984年（宋雍熙元年、西夏天授礼法延祚-53年；公元换算据两国年号对照）] \eventref{SYD-984-XIA-001}{夏州言，掩擊李繼遷，獲其母妻，俘千四百餘帳，繼遷走}（宋与西夏）。
  \item[985年（宋雍熙二年、西夏天授礼法延祚-52年；公元换算据两国年号对照）] \eventref{SYD-985-XIA-002}{辛丑，夏州行營破西蕃息利族，斬其代州刺史折羅遇并弟埋乞，又破保、洗兩族，降五十餘族}（宋与西夏）。
  \item[985年（宋雍熙二年、西夏天授礼法延祚-52年；公元换算据两国年号对照）] \eventref{SYD-985-XIA-053}{乙未，夏州李繼遷誘殺汝州團練使曹光實}（宋与西夏）。
\end{description}
\subsection{金}
\begin{description}
  \item[1115年（金收国元年、辽天庆五年、宋政和五年；干支、月日待核；公元换算据各方本纪年号对照）] \eventref{JIN-1115-EVENT-079}{完颜阿骨打称帝建金}（金）。
  \item[1118年] \eventref{JIN-1118-EVENT-080}{金与宋讨论共同攻辽}（金与宋）。
  \item[1118年] \eventref{NS-1118-SONG-JIN-CONTACT}{宋廷经海路和边境渠道接触金方，筹议夹攻辽朝}（宋与金）。
  \item[1119年] \eventref{NS-1119-EVENT-022}{宋金继续讨论夹攻辽朝与燕云处置}（宋与金）。
\end{description}
\subsection{南宋}
\begin{description}
  \item[1127年（宋建炎元年、金天会五年；干支、月日待核；公元换算据双方本纪年号对照）] \eventref{SS-1127-EVENT-111}{赵构即位为高宗，南宋朝廷开始重建}（南宋）。
  \item[1127年（南宋靖康二年、金天会5年；公元换算据两国年号对照）] \eventref{SYD-1127-JIN-009}{丁酉，金人立張邦昌為楚帝}（南宋与金）。
  \item[1127年（南宋靖康二年、金天会5年；公元换算据两国年号对照）] \eventref{SYD-1127-JIN-108}{二年春正月辛卯朔，命濟王栩、景王杞出賀金軍，金人亦遣使入賀}（南宋与金）。
  \item[1127年（南宋靖康二年、金天会5年；公元换算据两国年号对照）] \eventref{SYD-1127-JIN-162}{庚子，金人來取宗室，開封尹徐秉哲令民結保，毋藏匿}（南宋与金）。
\end{description}
\subsection{蒙古}
\begin{description}
  \item[1206年（元中统-53年；公元换算据《元史》本纪纪年）] \eventref{SYD-1206-MONGOL-001}{阿剌忽思即以是謀報帝，居無何，舉部來歸}（蒙古）。
  \item[1206年（元中统-53年；公元换算据《元史》本纪纪年）] \eventref{SYD-1206-MONGOL-053}{別里古台亦曰：「乃蠻欲奪我弧矢，是小我也，我輩義當同死}（蒙古）。
  \item[1206年（元中统-53年；公元换算据《元史》本纪纪年）] \eventref{SYD-1206-MONGOL-103}{不聽，各持馬乳橦疾鬭，奪忽兒真、火里真二哈敦以歸}（蒙古）。
  \item[1206年（元中统-53年；公元换算据《元史》本纪纪年）] \eventref{SYD-1206-MONGOL-150}{帝悅，曰：「以此眾戰，何憂不勝}（蒙古）。
\end{description}
\subsection{元及区域行动者}
\begin{description}
  \item[1271年（元至元八年；干支、月日待核；公元换算据《元史》本纪年号对照）] \eventref{YUAN-1271-EVENT-173}{忽必烈定国号元}（元）。
  \item[1272年（元至元9年；公元换算据《元史》本纪纪年）] \eventref{SYD-1272-YUAN-001}{:今遣汝等，水陸並進，布告遐邇，使咸知之}（元）。
  \item[1272年（元至元9年；公元换算据《元史》本纪纪年）] \eventref{SYD-1272-YUAN-145}{敕杖合丹，斥無入宿衞，謫往西川效死軍中，餘定罪有差}（元）。
  \item[1272年（元至元9年；公元换算据《元史》本纪纪年）] \eventref{SYD-1272-YUAN-217}{從其請，仍降璽書，遣使諭江陵府制置司及高達已下官吏軍民}（元）。
\end{description}
\section{问题索引}
\subsection{战争、城防与道路}
\begin{description}
  \item[辽与渤海] \eventref{LIAO-0926-EVENT-028}{辽军攻灭渤海，设置东丹等安排}（926年）。
  \item[辽与后晋] \eventref{LIAO-0947-EVENT-032}{辽军进入开封，后因补给与抵抗压力撤离}（947年）。
  \item[辽与宋] \eventref{LIAO-0979-EVENT-034}{辽军在高梁河击退北伐宋军}（979年）。
  \item[辽与宋] \eventref{LIAO-0986-EVENT-036}{辽在雍熙北伐中应对宋军多路进攻并取胜}（986年）。
  \item[后周禁军与赵匡胤集团] \eventref{NS-0960-CHENQIAO}{陈桥兵变后赵匡胤即帝位，改国号宋}（960年正月）。
  \item[宋] \eventref{NS-0961-MILITARY-POWER}{太祖以宴饮、任官和调动方式削弱禁军宿将的独立兵权}（961年前后）。
  \item[宋与荆南] \eventref{NS-0963-JINGNAN}{宋军入荆南，结束高氏政权}（963年）。
  \item[宋与后蜀] \eventref{NS-0964-LATER-SHU-CAMPAIGN}{宋军发动伐后蜀战役}（964年冬）。
\end{description}
\subsection{财政、物资与流通}
\begin{description}
  \item[辽] \eventref{LIAO-0938-EVENT-031}{辽以幽州为南京，经营燕云城市与交通节点}（938年）。
  \item[宋] \eventref{SYD-962-NS-003}{八月癸巳，蔡河務綱官王訓等四人坐以糠土雜軍糧，磔於市}（962年（宋建隆三年；公元换算据《宋史》本纪纪年））。
  \item[宋] \eventref{SYD-965-NS-173}{丙申，赦蜀，歸俘獲，除管内逋賦，免夏税及沿徵物色之半}（965年（宋乾德三年；公元换算据《宋史》本纪纪年））。
  \item[宋与辽] \eventref{SYD-970-LIAO-006}{癸亥，定州駐泊都監田欽祚敗契丹於遂城}（970年（宋開寶三年、辽保宁2年；公元换算据两国年号对照））。
  \item[宋] \eventref{SYD-972-NS-180}{丙午，遣使檢視水災田}（972年（宋開寶五年；公元换算据《宋史》本纪纪年））。
  \item[宋] \eventref{SYD-973-NS-181}{丙申，曹州饑，漕太倉米二萬石振之}（973年（宋開寶六年；公元换算据《宋史》本纪纪年））。
  \item[宋] \eventref{SYD-986-NS-194}{丙子，召曹彬、崔彥進、米信歸闕，命田重進屯定州，潘美還代州}（986年（宋雍熙三年；公元换算据《宋史》本纪纪年））。
  \item[宋] \eventref{SYD-994-NS-202}{貝州言驍捷卒劫庫兵為亂，推都虞候趙咸雍為帥，轉運使王嗣宗率屯兵擊敗之，擒咸雍，磔於市}（994年（宋淳化五年；公元换算据《宋史》本纪纪年））。
\end{description}
\subsection{继承、官制与地方权力}
\begin{description}
  \item[契丹] \eventref{LIAO-0907-EVENT-026}{耶律阿保机取得汗位，开始重组部族权力}（907年）。
  \item[契丹] \eventref{LIAO-0916-EVENT-027}{阿保机称帝并建国号契丹}（916年）。
  \item[辽] \eventref{LIAO-0927-EVENT-029}{耶律德光即位为辽太宗}（927年）。
  \item[辽与后晋] \eventref{LIAO-0936-EVENT-030}{辽支持石敬瑭建立后晋，十六州问题进入北方政治结构}（936年）。
  \item[辽] \eventref{LIAO-0969-EVENT-033}{辽景宗即位，萧绰逐步取得政治作用}（969年）。
  \item[辽] \eventref{LIAO-0982-EVENT-035}{辽圣宗即位，萧太后摄政}（982年）。
  \item[辽与宋] \eventref{LIAO-0986-EVENT-036}{辽在雍熙北伐中应对宋军多路进攻并取胜}（986年）。
  \item[宋] \eventref{NS-0961-MILITARY-POWER}{太祖以宴饮、任官和调动方式削弱禁军宿将的独立兵权}（961年前后）。
\end{description}
\subsection{盟约、使节与跨境关系}
\begin{description}
  \item[宋] \eventref{SYD-960-NS-168}{丙辰，南唐主李景、吳越王錢俶遣使以御服、錦綺、金帛來賀}（960年（宋建隆元年；公元换算据《宋史》本纪纪年））。
  \item[宋与辽] \eventref{SYD-963-LIAO-002}{己亥，契丹幽州岐溝關使柴廷翰等來降}（963年（宋乾德元年、辽应历13年；公元换算据两国年号对照））。
  \item[宋与辽] \eventref{SYD-963-LIAO-086}{戊寅，北漢引契丹兵攻平晉，遣洺州防禦使郭進等救之}（963年（宋乾德元年、辽应历13年；公元换算据两国年号对照））。
  \item[宋与辽] \eventref{SYD-966-LIAO-004}{丁巳，契丹天德軍節度使于延超與其子來降}（966年（宋乾德四年、辽应历16年；公元换算据两国年号对照））。
  \item[宋] \eventref{SYD-966-NS-007}{安國軍節度使羅彥瓌等敗北漢於靜陽，擒其將鹿英}（966年（宋乾德四年；公元换算据《宋史》本纪纪年））。
  \item[宋] \eventref{SYD-967-NS-175}{北漢鴻唐砦招收指揮使樊暉以砦來降}（967年（宋乾德五年；公元换算据《宋史》本纪纪年））。
  \item[宋与辽] \eventref{SYD-969-LIAO-087}{命彰德軍節度使韓仲贇為北面都部署，彰義軍節度使郭延義副之，以防契丹}（969年（宋開寶二年、辽保宁1年；公元换算据两国年号对照））。
  \item[宋] \eventref{SYD-969-NS-177}{賜宰相、樞密使、翰林學士、節度、觀察使襲衣金帶}（969年（宋開寶二年；公元换算据《宋史》本纪纪年））。
\end{description}
\section{材料与图件入口}
\begin{description}
  \item[正史与编年] \hyperref[app:sources]{“来源与史料说明”}说明《宋史》《辽史》《金史》《元史》《续资治通鉴长编》及《建炎以来系年要录》的不同用途。
  \item[文书、碑刻与考古] \hyperref[app:sources]{来源索引}列明黑水城文书、契丹／女真／蒙古等碑刻、宋元契约、城址和港口材料如何补正官修年序。
  \item[空间总览] \hyperref[fig:vol04:overview:multi-polity-spatial]{“宋、辽、西夏、金与元的多政权空间关系示意”}。
  \item[时间轴] \hyperref[fig:vol04:overview:multi-polity-chronology]{“宋辽夏金元并立年序示意”}。
  \item[北方与河西边地] \hyperref[fig:vol04:liao-western-xia:multi-polity-frontiers]{“辽、西夏与宋的多政权边地关系示意”}。
  \item[南宋海运] \hyperref[fig:vol04:southern-song:maritime-transport-network]{“南宋海域与运输网络示意”}。
  \item[蒙古征服与元治理] \hyperref[fig:vol04:mongol-conquest:empire-formation]{“蒙古征服与帝国形成的区域关系示意”}；\hyperref[fig:vol04:yuan:governance-provincial-network]{“元朝行省、交通与治理关系示意”}。
\end{description}
